% ============================================================
%  Глава 3. Оценка rainbow-опционов с помощью методов
%           машинного обучения
% ============================================================
\chapter{Оценка rainbow-опционов с помощью методов
         машинного обучения}
\label{ch:ml}

\section{Зачем использовать ML в оценке деривативов}
\label{sec:ml-why}

Для рассматриваемых в работе экзотических деривативов классический
аппарат оценки имеет ряд ограничений. Во-первых, для большинства
интересующих нас инструментов отсутствуют удобные замкнутые формулы и
основным методом становится метод Монте-Карло, требующий значительных
вычислительных затрат, особенно при необходимости пересчёта цен и
греков на большом числе сценариев в реальном времени. Во-вторых,
дифференциация цены по риск-факторам в Монте-Карло требует либо
конечных разностей (нестабильно при разумных размерах симуляций), либо
методов pathwise/likelihood ratio, реализация которых для сложных
инструментов трудоёмка.

Альтернативой является обучение нейронной сети как функции, отображающей
вектор риск-факторов в цену дериватива. При этом, благодаря
автоматическому дифференцированию, греки получаются как простой
дополнительный выход той же сети без отдельных пересчётов. Эта идея
систематически развивалась в литературе по
\emph{deep~hedging}~\cite{buhler2019deephedging} и обзорах применения
нейросетей для оценки и хеджирования
деривативов~\cite{horvath2021deepvol, joas2024}.

\section{Обучение остатка декомпозиции цены}
\label{sec:ml-residual}

Прямое обучение модели <<все факторы $\to$ цена rainbow-опциона>> сложно
по двум причинам:
(1)~функция отображения существенно \emph{нелинейна} и зависит от большого
числа коррелирующих переменных;
(2)~часть зависимости известна аналитически (через цены ванильных
опционов и форвардов на базовые активы), поэтому обучать её заново~---
расточительно.

Поэтому в работе используется следующий подход. Цена rainbow-опциона
раскладывается по~\eqref{eq:rainbow-decomp}, и отдельно оценивается
лишь \emph{остаточный} компонент $R$. Все аналитически считаемые
слагаемые вычисляются по замкнутым формулам Блэка--Шоулза и Маргрэйба
для спред-опциона, а остаток $R$ обучается на синтетических данных,
сгенерированных численно (см.~\cref{ch:data}).

Формально, целевая функция модели имеет вид
\begin{equation}
  R(x;\theta)
  \;=\; C_{\text{rb}}(x) - C_{\text{vanilla}}(x) - C_{\text{spread}}(x),
\end{equation}
где $x = (S^{(1)},S^{(2)},\sigma_1,\sigma_2,\rho,K,T,r)$ --- вектор
входных параметров, а $\theta$~--- веса модели. Эмпирически такой
подход даёт меньшую дисперсию ошибок и лучшую устойчивость по сравнению
с прямой оценкой, особенно <<вблизи>> страйка и при экстремальных
значениях корреляции.

\section{Прогнозирование цены за счёт прогноза риск-факторов
         \texorpdfstring{(\emph{for info})}{(for info)}}
\label{sec:ml-forecast}

Дополнительный, не входящий в основной эксперимент данной работы,
сюжет состоит в следующем. Раскладывая цену rainbow-опциона
по~\eqref{eq:rainbow-decomp}, можно прогнозировать каждое слагаемое в
отдельности и собирать прогноз цены инструмента как сумму прогнозов
своих компонент:

\begin{itemize}
  \item цены базовых активов~--- модели типа ARIMA на лог-доходностях;
  \item волатильности~--- модели стохастической волатильности
        (Heston, SABR) или GARCH-семейство;
  \item корреляции~--- модель Якоби~\eqref{eq:jacobi} или Фишера--OU
        \eqref{eq:fisher-ou};
  \item остаток~$R$~--- предобученная сетка, ср. \cref{sec:ml-residual}.
\end{itemize}

В рамках настоящей работы этот сценарий приведён лишь как замечание о
перспективе развития; основной эксперимент~--- сравнение моделей цены
на одной и той же временной отметке, а не прогноз будущей цены.
